%%%%%%%%%%%%%%%%%%%%%%%%%%%%%%%%%%%%%%%%%
% Dictionary
% LaTeX Template
% Version 1.0 (20/12/14)
%
% This template has been downloaded from:
% http://www.LaTeXTemplates.com
%
% Original author:
% Vel (vel@latextemplates.com) inspired by a template by Marc Lavaud
% Modified by:
% Natsuko Nakagawa
%
% License:
% CC BY-NC-SA 3.0 (http://creativecommons.org/licenses/by-nc-sa/3.0/)
%
%%%%%%%%%%%%%%%%%%%%%%%%%%%%%%%%%%%%%%%%%

\documentclass[a4paper,xelatex,ja=standard, enablejfam=true, jbase=8pt, nomag]{article}
%%%%%%%%%%%%%%%%%%%%%%%%%%%%%%%%%%%%%%%%%%%%%%%%%%%%
%%%             PACKAGES                         %%%
%%%%%%%%%%%%%%%%%%%%%%%%%%%%%%%%%%%%%%%%%%%%%%%%%%%%

% Document margins and spacings
\usepackage[margin=1.5cm,columnsep=5pt]{geometry} 

% Japanese font
\usepackage{fontspec}
\usepackage{zxjatype}
\usepackage[ipaex]{zxjafont}
\defaultfontfeatures{Ligatures=TeX}

% IPA font
\newfontfamily\ipafont{Charis SIL}
\newcommand\ipa[1]{{\ipafont #1}}

% Index
\usepackage{imakeidx}
\makeindex[columns=3, title=逆引き, intoc]

% Multiple columns
\usepackage{multicol}

\usepackage{stix,indentfirst}

% Headers & footers
\usepackage{fancyhdr} % Required for modifying headers and footers
\fancyhead[LE]{\textsf{\rightmark}} % Top left header, even page
\fancyhead[RO]{\textsf{\rightmark}} % Top right header, odd page
\fancyhead[LO]{} % Top left header, odd page
\fancyhead[RE]{} % Top right header, even page
\renewcommand{\headrulewidth}{1.4pt} % Rule under the header
\fancyfoot[C]{{\thepage}} % Bottom center footer
\pagestyle{fancy} % Use the custom headers and footers throughout the document

% 行間の調整
\renewcommand{\baselinestretch}{1.4}
\renewcommand{\arraystretch}{0.9}

% その他パッケージ
\usepackage{fancybox}

% 例文マーク
\newcommand{\reibun}{\ovalbox{\small 例}}

% 辞書エントリの定義
\newcommand{\entry}[5]{\hangindent = 1cm
\textbf{#1}\ \textbf{\ipa{#2}}{【#3】}\ {#4}{#5}
\markboth{#1}{#1}}
% \entry{仮名項目}{IPA}{品詞}{記述}{例文}

\newcommand{\ctd}[2]{\hspace{0.5cm}{#1}{#2}}
% \ctd{記述}{例文}

% Title, author(s), & date
\title{\Huge \bf タイトル}
\author{\huge \bf 著者名}
\date{}

%%%%%%%%%%%%%%%%%%%%%%%%%%%%%%%%%%%%%%%%%%%%%%%%%%%%
%%%             辞書本体部分                     %%%
%%%%%%%%%%%%%%%%%%%%%%%%%%%%%%%%%%%%%%%%%%%%%%%%%%%%

\begin{document}

\maketitle

\newpage

%%%%%%%%%%%%%%%%%%%%%%%%%%%%%%%%%%%%%%%%%%
%% エクセルテンプレート貼り付け開始     %%
%%%%%%%%%%%%%%%%%%%%%%%%%%%%%%%%%%%%%%%%%%

\section*{あ}\begin{multicols}{2}

\entry{あが ぴてぎ}{[\ipa{aga pitegi}] }{名}{赤畑（赤土の畑）。}{}\index{あかばたけ（あかつちのはたけ）}

\entry{あがあまじ}{[\ipa{aga-amadʒi}] }{名}{赤毛。}{}\index{あかげ}

\entry{あかき}{[\ipa{akaki}] }{名}{赤毛。}{}\index{あかげ}

\entry{あがじー}{[\ipa{agadʒii}] F}{名}{血。}{\reibun じー　んだん [\ipa{dʒii ndan}] (血が出た) 。}\index{ち}

\entry{あがじん}{[\ipa{aga-dʒin}] }{名}{赤い銭（十円玉など）。}{}\index{あかいぜに（じゅうえんだま）}

\entry{あがぬり}{[\ipa{aga-nuri}] }{名}{田んぼに入ったときに足につく赤い泥。}{\reibun あがぬり　かぺる　どー [\ipa{aga-nuri kḁp-e-ru doo}] (赤泥を被っているよ) 。}\index{あかいどろ}

\entry{あがぱな}{[\ipa{aga-pḁn̥a}] }{名}{赤花。}{}\index{あかばな}

\entry{あがはん}{[\ipa{agahan}] }{形}{赤い。}{}\index{あかい}

\entry{あがまみ}{[\ipa{agaˑmami}] }{名}{赤豆。}{}\index{あかまめ}

\entry{あがりぴる}{[\ipa{agari-pi̥r̥u}] }{名}{明け方。}{}\index{あけがた}

\entry{あき}{[\ipa{aki}] }{名}{秋。}{}\index{あき}

\entry{あぎさり}{[\ipa{agisari}] F}{名}{明け方日が昇る前の明るい時間帯。}{}\index{あけがた}

\entry{あぎとーふ}{[\ipa{agi-tooɸu}] }{名}{揚げ豆腐。}{}\index{あげどうふ}

\entry{あぎみっくゎ}{[\ipa{agi-mikkwa}] }{名}{あきめくら罵りことば。}{}\index{あきめくら}

\entry{あぐ}{[\ipa{agu}] }{名}{顎。}{}\index{あご}

\entry{あくび}{[\ipa{akubi}] F}{名}{あくび。}{}\index{あくび}

\entry{あごん}{[\ipa{agon}] F}{名}{アコウの木。}{}\index{あこう}

\entry{あこんき}{[\ipa{agon-ki}] }{名}{アコウの木。}{}\index{あこう}

\entry{あざ}{[\ipa{adza}] H}{名}{そば。}{\reibun あざご　くば [\ipa{adza=go ku-ba}] (そばに来い) 。}\index{そば}

\entry{あざ}{[\ipa{adza}] H}{名}{畑のあぜ道。}{\reibun ぴてぎぬ　あざ [\ipa{pitegi=nu adza}] (畑のあぜ道) 。}\index{はたけのあぜみち}

\entry{あし}{[\ipa{aʃi}] H}{名}{汗。}{\reibun あっつぁはったら　あし　ふき　よ [\ipa{attsa-ha-ttara aʃi ɸu̥ki jo}] (暑いので汗をかいた) 。}\index{あせ}

\entry{あすとぅ}{[\ipa{astu}] H}{名}{明後日。}{\reibun あすとぅ　ぴま　ねぬ？ [\ipa{astu pima nenu?}] (明後日暇はないか？) 。}\index{あさって}

\entry{あすまり}{[\ipa{asm̥arun}] }{動}{集まる。}{\reibun きゅーや　ぶらくぬ　そーかい　やりき　あすまり　ひより [\ipa{kjuuja burakunu sookai jariki asm̥ari hijori}] (今日は集落の総会なので、集まってください) 。}\index{あつまる}

\entry{あとぅむぬ}{[\ipa{atumunu}] F}{名}{胞衣。}{\reibun あとぅむぬ　んだん [\ipa{atumunu ndan}] (胞衣は出たか？) 。}\index{えな}

\entry{あばだり}{[\ipa{abadari}] F}{名}{裸。}{\reibun あばだり　しぇろ [\ipa{abadari ʃero}] (裸になっている) 。}\index{}

\entry{あばなび}{[\ipa{abanabi}] }{名}{揚げ物用の鍋。}{}\index{あげものようのなべ}

\entry{あばりしゃ}{[\ipa{abari-ʃa}] }{名}{(人が) 綺麗だ。}{\reibun うぬ　ぴとぅ　あばりしゃ　わら [\ipa{unu pi̥tu abari-ʃa wara}] (この人はきれいだね) 。うぬ　あっぱー　むがし　あばりしゃった　っちょ [\ipa{unu appaa mugaʃi abari-ʃa-tta ttʃo}] (このおばあさんは昔きれいだったらしい) 。}\index{きれいだ}

\entry{あふ}{[\ipa{aɸu}] F}{名}{灰汁（はいじる）灰を混ぜた水汚れを落とす作用がある。}{}\index{あく}

\entry{あぶく}{[\ipa{abuku}] H}{名}{泡。}{\reibun とふ　あくぬ　ふきけん [\ipa{toɸu abuku ɸu̥ki-ken}] (豆腐が泡を吹いてきた) 。}\index{あわ}

\entry{あまだりみじ}{[\ipa{amadari-midʒi}] }{名}{軒から落ちる雨水をためたもの。}{\reibun あまだりみじ　たみり [\ipa{amadari-midʒi tḁm̥ir-i}] (雨水をためろ) 。}\index{あまみず}

\entry{あみ}{[\ipa{ami}] F}{名}{飴。}{\reibun あみぬ　ながな　ふねぶどぅ　ぺれろ [\ipa{ami=nu naga=na ɸu̥n̥ebu=du per-e-ro}] (飴の中にみかんが入っている) 。}\index{あめ}

\entry{あみ}{[\ipa{ami}] H}{名}{雨。}{\reibun あみ　ぱらん [\ipa{ami pḁr̥-a-n}] (雨が晴れた (止んだ)) 。あみぬ　ながじゅどぅ　あらぎ　き [\ipa{ami=nu naga=ju=du arag-i k-i}] (雨の中を歩いてきた) 。}\index{あめ}

\entry{あみ}{[\ipa{ami}] F}{名}{飴。}{\reibun あみぬ　ながな　ふねぶどぅ　ぺれろ [\ipa{ami=nu naga=na ɸu̥n̥ebu=du per-e-ro}] (飴の中にみかんが入っている) 。}\index{あめ}

\end{multicols}\section*{い}\begin{multicols}{2}

\entry{いじ}{[\ipa{idʒi}] F}{名}{元気、気力。}{\reibun いじ　んだひよ [\ipa{idʒi ndahijo}] (頑張りなさい) 。}\index{げんき}

\end{multicols}\section*{う}\begin{multicols}{2}

\entry{うすな}{[\ipa{usɨ̥n̥a}] H}{名}{沖縄。}{\reibun うしなち　んぎた　くい [\ipa{usɨ̥na=tʃi ng-i-ta kui}] (沖縄に行ってこようね) 。}\index{おきなわ}

\entry{うび}{[\ipa{ubi}] F}{名}{いくら、どのぐらい。}{\reibun うびばがら　ぶさったら [\ipa{ubibagara busattara?}] (どれぐらい大きかった？) 。}\index{いくら、どのぐらい}

\entry{うり}{[\ipa{uri}] F}{名}{1それ。}{}\index{それ}

\ctd{2彼、彼女。}{}

\end{multicols}\section*{か}\begin{multicols}{2}

\entry{がば}{[\ipa{gaba}] F}{名}{垢。}{}\index{あか}

\entry{ごっかり}{[\ipa{gokkari}] }{名}{アカショウビン。}{}\index{あかしょうびん}

\entry{こんぎ}{[\ipa{koŋgi}] H}{名}{桑。}{\reibun こんぎんかんやー　ふくんどぅ　たかはろ [\ipa{koŋginkanjaa ɸu̥kundu takaharo}] (桑の木より福木が高い) 。}\index{くわ}

\entry{しぱり}{[\ipa{spari}] H}{名}{小便。}{}\index{しょうべん、にょう}

\entry{じん}{[\ipa{dʒin}] F}{名}{頂、てっぺん。}{\reibun やまぬ　じん [\ipa{jama=nu dʒin}] (山の頂上) 。}\index{いただき}

\entry{じん}{[\ipa{dʒin}] F}{名}{お盆、お膳。}{\reibun じん　んだひば [\ipa{dʒin ndah-i-ba}] (膳を出せ) 。}\index{おぼん}

\entry{すき}{[\ipa{ski}] H}{名}{鋤（すき）。}{\reibun だー　すき　からひ [\ipa{daa ski kḁr̥ahi}] (君の鋤（すき）を貸しなさい) 。すき　からひ　ひ　おらぬ？ [\ipa{ski kḁr̥ahi hi oranu}] (鋤（すき）を貸して下さいませんか？) 。}\index{すき}

\entry{っふぁ}{[\ipa{ɸɸa}] F}{名}{鞍。}{\reibun んまぬ　っふぁ [\ipa{mma=nu ɸɸa}] (馬の鞍) 。}\index{くら}

\entry{どぅー}{[\ipa{duu}] F}{名}{自分。}{\reibun どぅーぬ　ひがち　かえれん [\ipa{duunu higatʃi kaeren}] (自分の家に帰った) 。}\index{じぶん}

\entry{はこち}{[\ipa{hakotʃi}] }{名}{顎。}{}\index{あご}

\entry{ばった}{[\ipa{batta}] F}{名}{1腹。}{}\index{はら}

\ctd{2内臓、腸。}{}

\entry{ぱん}{[\ipa{pan}] H}{名}{歯。}{\reibun ふちぬ　ぱんどぅ　やーむー [\ipa{ɸu̥tʃi=nu pan=du jaamuu}] (歯が痛い) 。}\index{は}

\entry{ふちの あご}{[\ipa{ɸutʃino ago}] }{名}{顎。}{}\index{あご}

\entry{んに}{[\ipa{nni}] F}{名}{胸。}{}\index{むね}


%%%%%%%%%%%%%%%%%%%%%%%%%%%%%%%%%%%%%%%%%%
%% エクセルテンプレート貼り付け終わり   %%
%%%%%%%%%%%%%%%%%%%%%%%%%%%%%%%%%%%%%%%%%%

\end{multicols}

\printindex



\end{document}
